\documentclass{article}
\usepackage[cm]{fullpage}
\usepackage{layout}
\usepackage{hyperref}   % For links
\usepackage{array}      % For tables
\usepackage{graphicx} % Required for inserting images
\usepackage{xcolor}
\usepackage[dvipsnames]{xcolor}
\usepackage{listings}   % For coding
\usepackage{amssymb}    % For square
\usepackage{amsmath}    % For cases
\usepackage{multicol}   % For multiple columns

\title{Teoremas, Proposiciones y Corolarios}
\author{Bander}
\date{Agosto 2025}

\begin{document}
%\layout para ver el layout actual
\maketitle

\begin{flushleft}
  Este pdf contiene solamente los \textbf{ENUNCIADOS} de Teoremas, Proposiciones y Corolarios vistos
  en la parte teórica de la materia Complejidad Computacional (1C2025).
\end{flushleft}

\begin{flushleft}
  \LARGE \textbf{Teoremas importantes}
\end{flushleft}
\begin{itemize}
  \item \hyperlink{cooklevin}{\textcolor{Bittersweet}{Teorema de Cook Levin}}
  \item \hyperlink{jerTdet}{\textcolor{Bittersweet}{Jerarquía de Tiempos Determinísticos}}
  \item \hyperlink{jerTnodet}{\textcolor{Bittersweet}{Jerarquía de Tiempos no Determinísticos}}
  \item \hyperlink{ladner}{\textcolor{Bittersweet}{Teorema de Ladner}}
  \item \hyperlink{jerEspacio}{\textcolor{Bittersweet}{Jerarquía de Espacio}}
  \item \hyperlink{savitch}{\textcolor{Bittersweet}{Teorema de Savitch}}
  \item \hyperlink{immerman}{\textcolor{Bittersweet}{Teorema de Immerman-Szelepcsényi}}
  \item \hyperlink{baker}{\textcolor{Bittersweet}{Teorema de Baker, Gill, Solovay}}
  \item \hyperlink{karpLipton}{\textcolor{Bittersweet}{Teorema de Karp-Lipton}}
\end{itemize}

\pagebreak

\begin{flushleft}
  \textbf{Proposición 1:} 
  \textit{Sea $f : \{0,1\}^* \rightarrow \{0,1\}^*$, sea T una función construible en
  tiempo y sea $\Gamma$ un alfabeto. Si f es computable en tiempo $T(n)$ por una máquina $M = 
  (\Gamma, Q, \delta)$, entonces f es computable en tiempo $O(log|\Gamma| \cdot T(n))$ por 
  una máquina $M' = (\Sigma, Q', \delta')$ donde $\Sigma = \{0,1,\triangleright, \square \}$
  es el alfabeto estándar.}\\[0.5em]
\end{flushleft}

\begin{flushleft}
  \textbf{\hypertarget{prop2}{Proposición 2:}} 
  \textit{Sea $f : \{0,1\}^* \rightarrow \{0,1\}^*$ y sea $T$ una función construible en tiempo. 
  Si $f$ es computable en tiempo $T(n)$ por una máquina estándar de $k \geq 3$ cintas (entrada, 
  salida y $k - 2$ cintas de trabajo), entonces $f$ es computable en tiempo $O(T(n)^2)$ por una 
  máquina de cinta única.}\\[0.5em]
\end{flushleft}

\begin{flushleft}
  \textbf{Proposición 3:}
  \textit{Sea $f : \{0,1\}^* \rightarrow \{0,1\}^*$ y sea $T$ una función construible en tiempo. Si $f$
  es computable en tiempo $T(n)$ por una máquina estándar, entonces hay una máquina oblivious que 
  computa $f$ en tiempo $O(T(n)^2)$.}\\[0.5em]
\end{flushleft}

\begin{flushleft}
  \textbf{Proposición 4:}

  \textit{Sea $f : \{0,1\}^* \rightarrow \{0,1\}^*$ y $T$ una función construible en tiempo. Si $f$
  es computable por una máquina con cintas bi-infnitas en tiempo $T(n)$, entonces $f$ es computable
  por una máquina estándar en tiempo $O(T(n))$}\\[0.5em]
\end{flushleft}

\begin{flushleft}
  \textbf{Teorema 1:}
  \textit{[Turing 1936] halt no es computable.}\\[0.5em]
\end{flushleft}

\begin{flushleft}
  \textbf{\hypertarget{teo2}{Teorema 2:}}
  \textit{Existe una máquina U que computa la función $u(\langle i,x \rangle) = M_i(x)$. Más aún, 
  si $M_i$ con entrada $x$ termina en $t$ pasos, entonces $U$ con entrada $\langle i,x \rangle$ termina
  en $c \cdot t \cdot  log(t)$ pasos, donde $c$ depende solo de $i$.}
\end{flushleft}

\begin{flushleft}
  \textbf{Teorema 3:}
  \textit{Existe una máquina $\widetilde{U}$ que computa la función $\widetilde{u}(\langle i,t,x \rangle)$
  en tiempo $c \cdot t \cdot log(t)$, donde $c$ depende solo de $i$.} 
\end{flushleft}

\begin{flushleft}
  \textbf{Teorema 4:}
  $\mathsf{P} \subseteq \mathsf{NP}$ \\[0.5em]
\end{flushleft}

\begin{flushleft}
  \textbf{Teorema 5:}
  $\mathsf{NP} = \bigcup_{c \in \mathbb{N}} \mathsf{NDTime}(n^c)$
\end{flushleft}

\begin{flushleft}
  \textbf{Teorema 6:}
  \textit{Existe una máquina no-determinística $NU$ tal que $NU$ acepta $(\langle i,x \rangle)$ sii $N_i$
  acepta $x$ y si $N_i$ corre en tiempo $T(n)$ entonces $NU(\langle i,x \rangle)$ decide si $N_i$ acepta o 
  rechaza $x$ en $c \cdot T(|x|)$ pasos, donde $c$ depende solo de $i$.}
\end{flushleft}

\begin{flushleft}
  \textbf{Teorema 7:}
  \textit{La relación $\leq_p$ es transitiva.}
\end{flushleft}

\begin{flushleft}
  \textbf{\hypertarget{teo8}{Teorema 8:}}
  \textit{Si }$\mathsf{NP}$-hard $\cap \ \mathsf{P} \neq \emptyset$\textit{, entonces} $\mathsf{P} = 
  \mathsf{NP}$.
\end{flushleft}

\begin{flushleft}
  \textbf{Teorema 9:}
  \textit{Si }$\mathcal{L} \in \mathsf{NP}$-completo\textit{, entonces} $\mathcal{L} \in \mathsf{P} \iff 
  \mathsf{P} = \mathsf{NP}$.
\end{flushleft}

\begin{flushleft}
  \textbf{Proposición 6:}
  $\mathsf{TMSAT} \in \mathsf{NP}$-completo. \\[0.5em]

  \textbf{\textcolor{OliveGreen}{Recordatorio:}} \\[0.5em]

  $\mathsf{TMSAT} = \{ \langle y,x,1^n,t^t \rangle \ | \ \exists u \in \{ 0,1 \}^n \ M_y(xu) = 1$ en a lo 
  sumo $t$ pasos$\}$ 
\end{flushleft}

\begin{flushleft}
  \textbf{Proposición 7:}
  \textit{Sea $\varphi(p_1,\dots, p_n)$ una fórmula booleana. Si $v$ y $v'$ son dos valuaciones tales que
  $v(p_i) = v'(p_i)$ para $i \in \{ 1, \dots, n\}$, entonces $v \models \varphi(p_1,\dots, p_n)$ sii
  $v' \models \varphi(p_1,\dots, p_n)$.}
\end{flushleft}

\begin{flushleft}
  \textbf{\hypertarget{teo10}{Teorema 10:}}
  $\mathsf{SAT}$, $3\mathsf{SAT} \in \mathsf{NP}$ 
\end{flushleft}

\begin{flushleft}
  \textbf{\hypertarget{prop8}{Proposición 8:}}
  \textit{Para toda $F : \{ 0,1 \}^\ell \rightarrow \{ 0,1 \}$ existe una fórmula booleana $\varphi_F(p_1,
  \dots, p_\ell)$ en CNF tal que $v \models \varphi$ sii $F(v) = 1$ para todo $v \in \{ 0,1 \}^\ell$. Más
  aún, $\varphi_F$ se computa en tiempo polinomial a partir de $\langle F \rangle$ y tiene tamaño $O(\ell
  2^\ell)$.} 
\end{flushleft}

\begin{flushleft}
  \textbf{Corolario 1:}
  \textit{Para toda $F : \{ 0,1 \}^\ell \rightarrow \{ 0,1 \}^k$ existe una fórmula booleana $\varphi(
  p_1, \dots, p_{\ell + k})$ en CNF tal que $uv \models \varphi_F$ sii $F(u) = v$ para todo $u \in \{
  0,1 \}^\ell$, $v \in \{0,1 \}^k$. Más aún, $\varphi_F$ se computa en tiempo polinomial a partir de 
  $\langle F \rangle$ y tiene tamaño $O((\ell + k)2^{\ell + k})$.}
\end{flushleft}

\begin{flushleft}
  \textbf{Teorema 11: \hypertarget{cooklevin}{\textcolor{Bittersweet}{(Teorema de Cook Levin)}}}
  $\mathsf{SAT} \in \mathsf{NP}$-hard. 
\end{flushleft}

\begin{flushleft}
  \textbf{Corolario 2:}
  $\mathsf{SAT}$, $3\mathsf{SAT} \in \mathsf{NP}$-completos. 
\end{flushleft}

\begin{flushleft}
  \textbf{Proposición 9:}
  $\mathsf{INDSET} \in \mathsf{NP}$-completo.
\end{flushleft}

\begin{flushleft}
  \textbf{Proposición 10:}
  $\mathsf{CAMHAM} \in \mathsf{NP}$-completo.
\end{flushleft}

\begin{flushleft}
  \textbf{Teorema 12:}
  $\mathsf{TSP} \in \mathsf{NP}$-completo.
\end{flushleft}

\begin{flushleft}
  \textbf{Proposición 11:}
  $\mathsf{KNAPSACK} \in \mathsf{NP}$-completo
\end{flushleft}

\begin{flushleft}
  \textbf{Teorema 13:}
  \textit{Si} $\mathsf{P} = \mathsf{NP}$ \textit{entonces} $\mathsf{ExpTime} = \mathsf{NExpTime}$ 
\end{flushleft}

\begin{flushleft}
  \textbf{Teorema 14: \hypertarget{jerTdet}{\textcolor{Bittersweet}{(Jerarquía de Tiempos Determinísticos)}}}
  \textit{Si $f$, $g$ son construibles en tiempo y cumplen que $f(n) \cdot log(f(n)) = o(g(n))$ entonces $
  \mathsf{DTime}(f(n)) \subsetneq \mathsf{DTime}(g(n))$.} 
\end{flushleft}

\begin{flushleft}
  \textbf{Teorema 15: \hypertarget{jerTnodet}{\textcolor{Bittersweet}{(Jerarquía de Tiempos no Determinísticos)}}}
  \textit{Si $f$, $g$ son construibles en tiempo y cumplen que $f(n) \cdot log(f(n)) = o(g(n))$ entonces $
  \mathsf{NDTime}(f(n)) \subsetneq \mathsf{NDTime}(g(n))$.} 
\end{flushleft}

\begin{flushleft}
  \textbf{Teorema 16: \hypertarget{ladner}{\textcolor{Bittersweet}{(Teorema de Ladner)}}} 
  \textit{Si $\mathsf{P} \neq \mathsf{NP}$ entonces existe $\mathcal{L}$ tal que $\mathcal{L} \in \mathsf{NP}$,
  $\mathcal{L} \not\in \mathsf{P}$, y $\mathcal{L} \not\in \mathsf{NP}$}-completo. 
\end{flushleft}

\begin{flushleft}
  \textbf{Proposición 12:}
  $\mathsf{SAT}_H \in \mathsf{P} \Rightarrow \exists c \ \forall \ n H(n) \leq c$.
\end{flushleft}

\begin{flushleft}
  \textbf{Proposición 13:}
  $\exists c \ \exists^ {\infty} n \ H(n) \leq c \Rightarrow \mathsf{SAT}_H \in \mathsf{P}$.
\end{flushleft}

\begin{flushleft}
  \textbf{Proposición 14:}
  \textit{Existe $k$ tal que para toda fórmula booleana $\psi$, $|\psi| > k \Rightarrow |F_\psi| <
  |\psi|$.}
\end{flushleft}

\begin{flushleft}
  \textbf{Proposición 15:}
  $\mathsf{DTime}(S(n)) \subseteq \mathsf{Space}(S(n))$.
\end{flushleft}

\begin{flushleft}
  \textbf{Proposición 16:}
  $\mathsf{Space}(S(n)) \subseteq \mathsf{NSpace}(S(n))$.
\end{flushleft}

\begin{flushleft}
  \textbf{Proposición 17:}
  \textit{Sea $M$ una máquina (determinística o no determinística) que usa espacio $S(n)$ tal que
  $S(n) \geq log(n)$. Existe una constante $c$ tal que para todo $x \in \{ 0,1 \}^*$, cada vértice
  de $G_{M, x}$ se puede describir usando $c \cdot S(|x|)$ bits ($c$ depende de $M$). Luego, $G_{M,x}$
  tiene $2^{c \cdot S(|x|)}$ nodos.}
\end{flushleft}

\begin{flushleft}
  \textbf{Proposición 18:}
  \textit{Si $S$ es construible en espacio, $\mathsf{NSpace}(S(n)) \subseteq \mathsf{DTime}(2^{O(S(n))})$.}
\end{flushleft}

\begin{flushleft}
  \textbf{Proposición 19:}
  $\mathsf{NP} \subseteq \mathsf{PSpace}$.
\end{flushleft}

\begin{flushleft}
  \textbf{Teorema 17: \hypertarget{jerEspacio}{\textcolor{Bittersweet}{(Jerarquía de Espacio)}}}
  \textit{Si $f$, $g$ son construibles en espacio y cumplen que $f(n) = o(g(n))$ entonces:}
\end{flushleft}
\begin{center}
  $\mathsf{Space}(f(n)) \subsetneq \mathsf{Space}(g(n))$.
\end{center}

\begin{flushleft}
  \textbf{Teorema 18:}
  $\mathsf{TQBF} \in \mathsf{PSpace}$.
\end{flushleft}

\begin{flushleft}
  \textbf{Proposición 20:}
  \textit{En tiempo polinomial se puede transformar cualquier QBF generalizada $\psi (y_1, \dots, y_m)$ 
  en una QBF $\psi'(y_1, \dots, y_m)$ posiblemente con variables libres que es equivalente, es decir, tal 
  que para todo $v \in \{ 0,1 \}^m$}
\end{flushleft}
\begin{center}
  $v \models \psi(y_1, \dots, y_m) \iff v \models \psi'(y_1, \dots, y_m)$.
\end{center}

\begin{flushleft}
  \textbf{Proposición 21:}
  $\mathsf{CHECKQBF} \leq_p \mathsf{TQBF}$.
\end{flushleft}

\begin{flushleft}
  \textbf{Proposición 22:}
  \textit{Sea $M$ una máquina (determinística o no-determinística) que usa espacio $S(n) \geq log(n)$ y sea
  $c$ una constante tal que para todo $x \in \{0,1\}^*$ cualquier configuración $C$ de un cómputo de $M$ con
  entrada $x$ se representa con $|\langle C \rangle| = c \cdot S(|x|)$ bits. Dado $x \in \{0,1\}^*$, podemos
  computar en tiempo polinomial en $S(n)$ una fórmula $\varphi_{M,x}(\overline{s}, \overline{t})$ en CNF con
  variables libres $\overline{s} = s_1, \dots, s_{c \cdot S(|x|)}$, $\overline{t} = t_1, \dots, t_{c \cdot
  S(|x|)}$ tal que $\langle C \rangle \langle C' \rangle \models \varphi_{M,x}(\overline{s}, \overline{t})$
  sii hay una flecha de $C$ a $C'$ en $G_{M,x}$.}
\end{flushleft}

\begin{flushleft}
  \textbf{Teorema 19:}
  $\mathsf{TQBF} \in \mathsf{PSpace}$-hard.
\end{flushleft}

\begin{flushleft}
  \textbf{Teorema 20: \hypertarget{savitch}{\textcolor{Bittersweet}{(Teorema de Savitch)}}}
  $\mathsf{NSpace}(S(n)) \subseteq \mathsf{Space}(S(n)^2)$\textit{, de modo que} $\mathsf{PSpace} = \mathsf{
  NPSpace}$.
\end{flushleft}

\begin{flushleft}
  \textbf{Proposición 23:}
  $\mathsf{PATH} \in \mathsf{NL}$
\end{flushleft}

\begin{flushleft}
  \textbf{Proposición 24:}
  \textit{Si $\mathcal{L} \not\in \{\{0,1\}^*, \emptyset\}$ entonces $\mathcal{L}$ es $\mathsf{NL}$}-hard
  \textit{con respecto a $\leq_p$.}
\end{flushleft}

\begin{flushleft}
  \textbf{Proposición 25:}
  \textit{Sean $f,g$ computables implícitamente en $\mathsf{L}$. Entonces $g \circ f$ es computable 
  implícitamente en $\mathsf{L}$.}
\end{flushleft}

\begin{flushleft}
  \textbf{Teorema 21:}
  $\mathsf{PATH} \in \mathsf{NL}$-completo \textit{y por lo tanto} $\overline{\mathsf{PATH}} \in \mathsf{coNL}$-completo.
\end{flushleft}

\begin{flushleft}
  \textbf{Teorema 22:}
  \textit{$\mathcal{L} \in \mathsf{NL}$ sii existe un polinomio $p : \mathbb{N} \rightarrow \mathbb{N}$ y una
  máquina determinística $M$ con una cinta de entrada adicional de lectura de una única vez (lee y pasa a la 
  siguiente celda a la derecha pero no puede volver atrás) tal que:}
\end{flushleft}
\begin{enumerate}
  \item \textit{Para todo $x \in {0, 1}^*$, $x \in \mathcal{L}$ sii existe $u \in {0, 1}^{p(|x|)}$ tal que $M (x, u) = 
  1$; aquí $M (x, u)$ denota la salida de $M$ cuando la cinta de entrada tiene $x$ y la cinta adicional de lectura 
  de una única vez tiene $u$}
  \item \textit{$M$ usa espacio $O(log(n))$; en este modelo de máquina con cinta de entrada adicional, la cinta de entrada 
  y la cinta adicional no cuentan en el espacio (solo cuentan sus cintas de trabajo y salida).}
\end{enumerate}

\begin{flushleft}
  \textbf{Teorema 23: \hypertarget{immerman}{\textcolor{Bittersweet}{(Teorema de Immerman-Szelepcsényi)}}}
  $\overline{\mathsf{PATH}} \in \mathsf{NL}$.
\end{flushleft}

\begin{flushleft}
  \textbf{Corolario 4:}
  $\mathsf{NL} = \mathsf{coNL}$.
\end{flushleft}

\begin{flushleft}
  \textbf{Teorema 24:}
  \textit{Si $S(n) \geq log(n)$ es construible en espacio entonces $\mathsf{NSpace}(S(n)) = \mathsf{coNSpace}
  (S(n))$.}
\end{flushleft}

\begin{flushleft}
  \textbf{Proposición 26:}
  $\textsf{MAXINDSET} \in \Sigma ^p _2$.
\end{flushleft}

\begin{flushleft}
  \textbf{Proposición 27:}
  \textit{La jerarquía polinomial tiene las siguientes propiedades:}
\end{flushleft}
\begin{multicols}{4}
  \begin{itemize}
    \item $\Sigma^p_1 = \mathsf{NP}$
    \item $\Sigma^p_i \subseteq \Sigma^p_{i+1}$
    \item $\Pi^p_i \subseteq \Sigma^p_{i+1}$
    \item $\Pi^p_1 = \mathsf{coNP}$
    \item $\Sigma^p_i \subseteq \Pi^p_{i+1}$
    \item $\Pi^p_i \subseteq \Pi^p_{i+1}$
    \item $\mathsf{PH} = \bigcup_{i \geq 0} \Sigma^p_{i}$
  \end{itemize}
\end{multicols}
\begin{flushleft}
  La pregunta $\mathsf{P} \stackrel{?}{=} \mathsf{NP}$ se puede generalizar a $\Sigma^p_i \stackrel{?}{=} 
  \Sigma^p_{i+1}$; ninguna se conoce. Decimos que la jerarquía polinomial colapsa si $\mathsf{P} = \mathsf{PH}$.
\end{flushleft}

\begin{flushleft}
  \textbf{Proposición 28:}
  \textit{Si $\mathsf{P} = \mathsf{NP}$ entonces $\mathsf{PH} = \mathsf{P}$.}
\end{flushleft}

\begin{flushleft}
  \textbf{Proposición 29:}
  \textit{Las clases $\Sigma^p_i$ y $\Pi^p_i$ están cerradas hacia abajo por $\leq_p$. Es decir, para $
  \mathcal{L}' \leq_p \mathcal{L}$ tenemos que si $\mathcal{L} \in \Sigma^p_i$ entonces $\mathcal{L}' \in 
  \Sigma^p_i$ y si $\mathcal{L} \in \Pi^p_i$ entonces $\mathcal{L}' \in \Pi^p_i$.}
\end{flushleft}

\begin{flushleft}
  \textbf{Proposición 30:}
  \textit{Si existe $\mathcal{L} \in \mathsf{PH}$}-completo \textit{entonces existe $i$ tal que $\mathsf{PH}
  = \Sigma^p_i$.}
\end{flushleft}

\begin{flushleft}
  \textbf{Corolario 5:}
  \textit{Si la jerarquía polinomial no colapsa en ningún nivel entonces $\mathsf{PH} \neq \mathsf{PSpace}$.}
\end{flushleft}

\begin{flushleft}
  \textbf{Proposición 31:}
  \textit{Para todo $i > 0$, $\Sigma_i \mathsf{SAT} \in \Sigma^p_i$.}
\end{flushleft}

\begin{flushleft}
  \textbf{Proposición 32:}
  \textit{Para todo $i > 0$, $\Sigma_i \mathsf{SAT} \in \Sigma^p_i$}-hard.
\end{flushleft}

\begin{flushleft}
  \textbf{Corolario 6:}
  \textit{Para todo $i > 0$, $\Sigma_i \mathsf{SAT} \in \Sigma^p_i$}-completo.
\end{flushleft}

\begin{flushleft}
  \textbf{Proposición 33:}
  \textit{Para todo $i > 0$, $\Pi_i \mathsf{SAT} \in \Pi^p_i$}-completo.
\end{flushleft}

\begin{flushleft}
  \textbf{Proposición 34:}
  \textit{Si $\mathcal{X} \in \mathsf{P}$, entonces $\mathsf{P} = \mathsf{P}^\mathcal{X}$.}
\end{flushleft}

\begin{flushleft}
  \textbf{Proposición 35:}
  $\mathsf{ExpTime \subseteq P^{EXPCOM}}$.
\end{flushleft}

\begin{flushleft}
  \textbf{Proposición 36:}
  $\mathsf{NP^{EXPCOM} \subseteq ExpTime}$.
\end{flushleft}

\begin{flushleft}
  \textbf{Corolario 7:}
  $\mathsf{P^{EXPCOM} = NP^{EXPCOM}}$.
\end{flushleft}

\begin{flushleft}
  \textbf{Teorema 25: \hypertarget{baker}{\textcolor{Bittersweet}{(Teorema de Baker, Gill, Solovay)}}}
  \textit{Existen oráculos $\mathcal{A}$ y $\mathcal{B}$ tal que $\mathsf{P}^\mathcal{A} = \mathsf{NP}^
  \mathcal{A}$ y $\mathsf{P}^\mathcal{B} \neq \mathsf{NP}^\mathcal{B}$.}
\end{flushleft}

\begin{flushleft}
  \textbf{Teorema 26:}
  \textit{Para $i \geq 1$, $\Sigma^p_{i+1} = \mathsf{NP}^{\Sigma_i \mathsf{SAT}}$.}
\end{flushleft}

\begin{flushleft}
  \textbf{Teorema 27:}
  $\mathsf{P \subseteq P_{/poly}}$.
\end{flushleft}

\begin{flushleft}
  \textbf{Teorema 28:}
  $\mathsf{P \nsupseteq P_{/poly}}$.
\end{flushleft}

\begin{flushleft}
  \textbf{Proposición 37:}
  \textit{Sea $(C_n)_{n \in \mathbb{N}}$ una familia de circuitos de tamaño $S(n)$ tal que\footnote{Notamos
  $C_n(\varphi)$ en ve de $C_n(\langle \varphi \rangle)$ y lo mismo para $C'_n$} $C_n(\varphi) = \mathcal{X}
  \mathsf{SAT}(\varphi)$ para toda fórmula booleana $\varphi = \varphi(x_1, \dots , x_n)$ con $|\varphi| = n$.
  Entonces existe una familia de circuitos $(C'_n)_{n \in \mathbb{N}}$ de tamaño polinomial en $S(n)$ tal que
  para todo $n: C'_n$ tiene $n$ salidas y para toda fórmula booleana $\varphi = \varphi(x_1, \dots, x_n)$ 
  con $|\varphi| = n$, tenemos $\varphi(C'_n(\varphi)) = \mathcal{X}\mathsf{SAT}(\varphi)$.}
\end{flushleft}

\begin{flushleft}
  \textbf{Teorema 29: \hypertarget{karpLipton}{\textcolor{Bittersweet}{(Teorema de Karp-Lipton)}}}
  \textit{Si $\mathsf{NP \subseteq P_{/poly}}$, entonces $\mathsf{PH} = \Sigma^p_2$.}
\end{flushleft}

\begin{flushleft}
  \textbf{Proposición 38:}
  $\mathsf{minSize}(f) = O(n2^n)$ \textit{para toda} $f : \{0,1\}^n \rightarrow \{0,1\}$.
\end{flushleft}

\begin{flushleft}
  \textbf{Teorema 30:}
  \textit{Para todo $n > 1$, existe $f : \{0,1\}^n \rightarrow \{0,1\}$ tal que $minSize(f) > 2^n / 4n$.}
\end{flushleft}

\begin{flushleft}
  \textbf{Proposición 39:}
  $\mathsf{PARITY \in NC^1_{nu}}$.
\end{flushleft}

\begin{flushleft}
  \textbf{Proposición 40:}
  $\mathsf{PARITY \not\in AC^0_{nu}}$.
\end{flushleft}

\begin{flushleft}
  \textbf{Proposición 41:}
  $\mathsf{SUMA \in AC^0_{nu}}$.
\end{flushleft}

\begin{flushleft}
  \textbf{Proposición 42:}
  $\mathsf{NC^d_{nu} \subseteq AC^d_{nu} \subseteq NC^{d+1}_{nu}}$\textit{. Por lo tanto,} $\mathsf{AC_{nu} = 
  NC_{nu}}$.
\end{flushleft}

\begin{flushleft}
  \textbf{Teorema 31:}
  \textit{$\mathcal{L}$ es decidible por una familia de circuitos $\mathsf{P}$}-uniforme \textit{sii} $
  \mathcal{L} \in \mathsf{P}$.
\end{flushleft}

\begin{flushleft}
  \textbf{Proposición 43:}
  $\mathsf{PARITY \in NC^1}$.
\end{flushleft}

\begin{flushleft}
  \textbf{Proposición 44:}
  $\mathsf{SUMA \in AC^0}$.
\end{flushleft}

\begin{flushleft}
  \textbf{Proposición 45:}
  $\mathsf{NC^d \subseteq AC^d \subseteq NC^{d+1}}$\textit{. Por lo tanto,} $\mathsf{AC = NC}$.
\end{flushleft}

\begin{flushleft}
  \textbf{Proposición 46:}
  $\mathsf{NC \subseteq P}$.
\end{flushleft}

\begin{flushleft}
  \textbf{Teorema 32:}
  \textit{$\mathcal{L}$ tiene una solución paralela eficiente sii $\mathcal{L} \in \mathsf{NC}$.}
\end{flushleft}

\begin{flushleft}
  \textbf{Teorema 33:}
  $\mathsf{NC^1 \subseteq L}$.
\end{flushleft}

\begin{flushleft}
  \textbf{Teorema 34:}
  \textit{Sea $\mathsf{L \in NSpace}(S(n))$. Existe una familia de circuitos $(C_n)_{n \in \mathbb{N}}$ 
  con compuertas $\land$ y $\lor$ de fan-in arbitrario y una máquina determinística $M$ tal que para todo
  $n \geq 1$ y $x \in \{0,1\}^n$ tenemos que $x \in L$ sii $C_n(x) = 1$, $M (1^n) = \langle C_n \rangle$,
  $|C_n| = 2^{O(S(n))}$ y $prof(C_n) = O(S(n))$.}
\end{flushleft}

\begin{flushleft}
  \textbf{Corolario 8:}
  $\mathsf{NL \subseteq AC^1}$.
\end{flushleft}

\begin{flushleft}
  \textbf{Proposición 47:}
  $\mathsf{CIRC}$-$\mathsf{EVAL}$ \textit{es} $\mathsf{P}$-completo.
\end{flushleft}

\begin{flushleft}
  \textbf{Teorema 35:}
  \textit{Supongamos que $\mathcal{L}$ es $\mathsf{P}$}-completo:
\end{flushleft}
\begin{enumerate}
  \item $\mathcal{L} \in \mathsf{NC \iff P = NC}$.
  \item $\mathcal{L} \in \mathsf{L \iff P = L}$.
\end{enumerate}
\end{document}
